% Источники: exam/materials/преподаватель/2020/23-03-2020-Методические указания Домрачева Моделирование.doc, лабораторная работа №2; exam/materials/моделирование.pdf, раздел 48; GitHub HumsterProgrammer/iu9-notes/8sem/Моделирование/Моделирование_лаб-00_12-02-26.md; GitHub HumsterProgrammer/iu9-notes/8sem/Моделирование/Моделирование_лаб-00_09-02-26.md.
\section{Модель Галилея и модель Ньютона.}

\textbf{Динамическая модель} отражает изменение состояния исследуемой
системы или объекта в заданный или формируемый период времени. Примером
динамической модели является баллистическая модель.

\textbf{Баллистическая модель} -- есть решение баллистической задачи в
следующей постановке: моделировать движение артиллерийского снаряда.
Выстрел произведен с начальной скоростью $v_0$ (м/с), под углом к горизонту
$\alpha$; снаряд изготовлен из известного материала, имеет форму шара с
радиусом $r$ (м). Требуется построить траекторию полета снаряда, указать
точку падения снаряда и время полета.

Пусть поверхность Земли плоская, а ускорение свободного падения -- постоянная
$g = 9{,}81$ м/с$^2$. Согласно законам механики, при сделанных
предположениях движение тела в горизонтальном направлении является
равномерным, а в вертикальном -- равномерным с ускорением, равным $g$.

На снаряд действуют сила тяжести $F_T$ и сила лобового сопротивления $F_C$.
В силу второго закона Ньютона формируется математическая модель движения
артиллерийского снаряда. Предложенное решение носит название
\textbf{модели Ньютона}. В ряде случаев, в том числе для проверки адекватности
модели Ньютона, применяют упрощенную модель, не учитывающую сопротивление
воздуха, -- \textbf{модель Галилея}.

\subsection*{Модель Галилея}

Для модели Галилея:
\[
\begin{cases}
  x = (v_0 \cos \alpha)t,\\
  y = (v_0 \sin \alpha)t - \dfrac{gt^2}{2}.
\end{cases}
\]
Отсюда:
\[
\begin{cases}
  t = \dfrac{x}{v_0\cos\alpha},\\
  y = \dfrac{v_0\sin\alpha}{v_0\cos\alpha}x
      - \dfrac{gx^2}{2v_0^2\cos^2\alpha}
    = (\tan\alpha)x - \dfrac{g}{2v_0^2\cos^2\alpha}x^2.
\end{cases}
\]

\subsection*{Модель Ньютона}

Для модели Ньютона используется запись
\[
  m\overline{a}=m\overline{g}-\beta |\overline{v}|\,\overline{v}.
\]
В обозначениях лабораторного конспекта:
\[
\begin{cases}
  m\dfrac{du}{dt}=-\beta u\sqrt{u^2+w^2}, & \dfrac{dx}{dt}=u,\\
  m\dfrac{dw}{dt}=-mg-\beta w\sqrt{u^2+w^2}, & \dfrac{dy}{dt}=w,
\end{cases}
\]
где $x,y$ — координаты тела, $u,w$ — составляющие скорости в проекциях на оси $X,Y$, $\beta$ — коэффициент сопротивления воздуха.

\textbf{Алгоритм решения задачи по методическим указаниям:}
\begin{enumerate}
  \item Изучить постановку баллистической задачи и модель Ньютона.
  \item Привести графическую интерпретацию постановки задачи: выбрать начало
  координат, именовать оси, указать единицы измерений величин по каждой оси,
  указать все входные данные и силы, действующие на материальную точку.
  \item Построить в соответствии с моделью Ньютона систему дифференциальных
  уравнений первого порядка.
  \item Реализовать решение системы методом Рунге--Кутты.
  \item Оценить траекторию полета по предложенной модели, вывести
  апостериорные оценки относительной погрешности значений высот в заданной
  точке, провести анализ результатов.
  \item Определить местоположение границы (стены) и решить обратную задачу с
  целью определения минимальной начальной скорости, требуемой для перелета
  груза через преграду.
  \item Сформулировать вычислительные задачи, обеспечивающие решение задачи;
  указать их вид и тип.
  \item Провести расчет по моделям Ньютона и Галилея, сравнить полученные
  результаты.
\end{enumerate}

\clearpage
